% IDPFlow NSF Report --- Times 12pt, single spacing, 1-inch margins
% Project Summary: 1 page | Project Description: 6 pages | References separate
\documentclass[12pt]{article}

\usepackage[T1]{fontenc}
\usepackage{mathptmx}          % Times New Roman
\usepackage[margin=1in]{geometry}
\usepackage{setspace}
\singlespacing
\usepackage[parfill]{parskip}
\setlength{\parskip}{6pt}
\usepackage{titlesec}
\titleformat{\section}{\normalfont\bfseries\large}{}{0em}{}[\titlerule]
\titleformat{\subsection}{\normalfont\bfseries}{}{0em}{}
\titlespacing*{\section}{0pt}{14pt}{4pt}
\titlespacing*{\subsection}{0pt}{10pt}{2pt}

\usepackage{hyperref}
\hypersetup{colorlinks=true, linkcolor=blue, citecolor=blue, urlcolor=blue}
\usepackage[numbers,sort&compress]{natbib}
\usepackage{booktabs}
\usepackage{microtype}
\usepackage{amsmath}
\usepackage{amssymb}

% No header/footer except page numbers
\pagestyle{plain}

\begin{document}

% ============================================================
%  PROJECT SUMMARY  (exactly 1 page)
% ============================================================
\begin{center}
  {\large\bfseries Project Summary}\\[4pt]
  \textbf{IDPFlow: Sequence-Conditioned Flow Matching for\\
  Intrinsically Disordered Protein Conformational Ensembles}
\end{center}

\noindent\textbf{Overview.}
Intrinsically disordered proteins (IDPs) and disordered regions (IDRs) constitute 30--40\% of the
human proteome and are central to cancer, neurodegeneration, and a wide range of other diseases.
Yet unlike folded proteins, IDPs lack a stable three-dimensional structure --- they sample a broad,
flat conformational energy landscape --- making them resistant to structure-based computational
methods. The current state-of-the-art for fast ensemble generation, BBFlow \cite{wolf2025bbflow},
achieves orders-of-magnitude speedups over molecular dynamics (MD) for folded proteins by
conditioning a geometric deep learning model on the equilibrium backbone structure
$\mathbf{x}_\mathrm{eq}$. For IDPs, no such equilibrium structure exists, and the method fails.
This proposal introduces \textbf{IDPFlow}, the first sequence-conditioned SE(3) flow matching
model for IDP conformational ensembles, requiring no structural input whatsoever.

\noindent\textbf{Intellectual Merit.}
IDPFlow will demonstrate that protein language model embeddings (ESM2 \cite{lin2023esm2})
combined with predicted secondary structure propensities
provide a sufficient conditioning signal to learn accurate IDP conformational distributions.
The model will be trained on the IDRome dataset \cite{tesei2024idrome} --- 28,058
coarse-grained CALVADOS2 \cite{lohr2022calvados2} MD trajectories of human IDPs --- providing
unprecedented scale for IDP machine learning.
Three specific aims will be pursued: (1) design and train the IDPFlow architecture with a
sequence-based conditional prior on the SE(3)$^N$ manifold; (2) validate against experimental
IDP ensembles from the Protein Ensemble Database (PED) \cite{lazar2021ped} using
radius-of-gyration distributions, per-residue RMSF, contact frequency maps, and pairwise RMSD;
(3) apply IDPFlow to the disease-relevant p53 N-terminal transactivation domain (TAD) and
$\alpha$-synuclein to identify transient structural features and cryptic binding sites.
Preliminary re-analysis of a 200-IDP subset of IDRome shows that sequence-derived
physicochemical descriptors alone predict mean $R_g$ with $R^2 = 0.87$ (cross-validated),
supporting the feasibility of sequence-conditioned ensemble learning.

\noindent\textbf{Broader Impacts.}
All trained model weights and code will be released as open-source software (MIT license),
enabling the community to generate IDP ensembles from sequence in under one second.
A publicly searchable database of ensemble predictions for over 10,000 human IDRs will be
deposited, expanding the utility of the IDRome resource to users who lack computational
infrastructure for MD simulation.
Two peer-reviewed publications are planned, targeting a machine learning venue (NeurIPS or
ICLR) and a structural biology journal (\textit{Nature Structural \& Molecular Biology} or
\textit{eLife}).
The disease application component will deliver cryptic binding site predictions for 20
disease-associated IDPs as a public resource for medicinal chemistry campaigns targeting
currently undruggable proteins.
Training of graduate and undergraduate researchers in the intersection of deep learning and
structural biology is an integral component of the project.

\noindent\textbf{Keywords:} intrinsically disordered proteins, conformational ensembles,
flow matching, SE(3) equivariance, protein language models, drug discovery, p53, $\alpha$-synuclein.

\newpage

% ============================================================
%  PROJECT DESCRIPTION  (6 pages)
% ============================================================
\begin{center}
  {\large\bfseries Project Description}
\end{center}

% -----------------------------------------------------------
\section{Overview and Specific Aims}
% -----------------------------------------------------------

Proteins are the molecular machines of life, and their functions depend not only on structure
but on dynamics. At physiological temperature, a protein samples an ensemble of conformations
according to the Boltzmann distribution,
$p(\mathbf{x}) \propto \exp(-E(\mathbf{x}) / k_B T)$.
For most globular proteins, this ensemble is tightly clustered around a single low-energy
equilibrium conformation that is accessible to X-ray crystallography or cryo-EM.
However, approximately 30--40\% of the human proteome contains intrinsically disordered
regions (IDRs) that lack a stable equilibrium structure
\cite{uversky2014idp_review,dyson2005idp}.
These regions are not merely flexible loops; they are fully functional molecular entities that
engage in protein--protein and protein--nucleic-acid interactions, mediate post-translational
modification cascades, and drive liquid--liquid phase separation.
Importantly, they are over-represented among disease-associated and clinically validated drug
targets: the p53 tumor suppressor, $\alpha$-synuclein in Parkinson's disease, and Tau in
Alzheimer's disease are all predominantly disordered \cite{uversky2008idp_disease}.

Generating accurate conformational ensembles for IDPs is therefore of major biological and
medical importance, yet it remains unsolved computationally.
MD simulation is the gold standard for ensemble accuracy, but IDPs require microsecond-to-millisecond
timescales to adequately sample conformational space --- 100- to 1,000-fold more than folded
proteins \cite{mittag2007idp_nmr}.
Recent geometric deep learning models, notably BBFlow \cite{wolf2025bbflow} and
AlphaFlow \cite{jing2024alphaflow}, achieve orders-of-magnitude speedups for folded proteins
by conditioning on an equilibrium backbone structure $\mathbf{x}_\mathrm{eq}$.
This conditioning is undefined for IDPs, and preliminary experiments show these models fail
or produce highly biased ensembles when forced to operate on AlphaFold2 predictions of
disordered regions, where predicted local distance difference test (pLDDT) scores are low
and the structural prediction is physically meaningless \cite{ruff2021af2idp}.

\noindent\textbf{Central hypothesis.}
Sequence-derived features --- specifically pretrained protein language model embeddings (ESM2
\cite{lin2023esm2}), predicted secondary structure propensities, and physicochemical
descriptors capturing charge patterning and hydrophobicity --- constitute a sufficient
conditioning signal to learn accurate IDP conformational distributions within an SE(3) flow
matching framework, entirely replacing the need for a structural input.

\noindent\textbf{Specific Aims:}

\noindent\textbf{Aim 1 --- IDPFlow architecture and training.}
We will design and implement IDPFlow, an SE(3) flow matching model conditioned on a
learned sequence representation $\boldsymbol{\phi}_\mathrm{seq}$ derived from ESM2
embeddings, predicted secondary structure propensities, and physicochemical descriptors.
The model will be trained on 23,000 IDP CALVADOS2 trajectories from the IDRome dataset
\cite{tesei2024idrome}.

\noindent\textbf{Aim 2 --- Validation against experimental IDP ensembles.}
IDPFlow will be benchmarked against experimental ensembles from the PED database
\cite{lazar2021ped} and against a held-out set of 2,558 IDRome trajectories, using
$R_g$ Wasserstein distance, per-residue RMSF Pearson correlation, contact frequency map
Frobenius distance, and pairwise C$_\alpha$ RMSD distribution MAE as quantitative metrics.

\noindent\textbf{Aim 3 --- Disease applications.}
IDPFlow will be applied to the p53 N-terminal transactivation domain (TAD, residues 1--73)
and full-length $\alpha$-synuclein (140 residues) to predict conformational ensembles and
identify cryptic binding sites using ensemble-based pocket detection.

% -----------------------------------------------------------
\section{Significance}
% -----------------------------------------------------------

\subsection{Background: intrinsically disordered proteins and their biology}

The classical structure--function paradigm --- one sequence, one fold, one function --- accounts
for the majority of biochemistry textbook examples but misrepresents a large fraction of
biologically critical proteins \cite{dyson2005idp}.
IDPs exhibit highly heterogeneous conformational behavior ranging from fully extended
random-coil states to transient secondary structure elements that form when binding partners
are present (folding-upon-binding), or even undergo complete folding into stable tertiary
structures in a partner-induced manner \cite{wright2009idp_binding}.
The diversity of IDP conformational behavior is the direct product of their sequence
composition: IDPs tend to have low mean hydrophobicity, high net charge, and low sequence
complexity \cite{uversky2014idp_review}, properties that prevent the hydrophobic collapse
driving the folding of globular proteins.

The biological importance of IDPs spans virtually every major disease area.
The p53 tumor suppressor, which is mutated or functionally inactivated in approximately
50\% of all human cancers, has a fully disordered N-terminal TAD that forms a transient
$\alpha$-helix at residues 19--26 (Box~1) upon binding to the MDM2 oncoprotein
\cite{raj2022p53}.
The interaction between this transient helix and MDM2 is the mechanistic basis for p53
inactivation in many cancers, and pharmacological disruption of this interface --- by
stabilizing alternative IDP conformations or by occupying the MDM2 binding groove --- is an
active area of drug discovery.
$\alpha$-Synuclein, a 140-residue IDP that is central to Parkinson's disease pathology,
populates an ensemble of disordered conformations in which transient long-range contacts
between the N-terminus and the aggregation-prone NAC domain (residues 61--95) play an
inhibitory role in fibril nucleation \cite{bhatt2021asyn}.
Tau, TDP-43, and FUS are further examples of disease-relevant IDPs for which ensemble
characterization is a prerequisite for mechanistic understanding \cite{uversky2008idp_disease}.

\subsection{Current computational methods and their limitations}

MD simulation is the established tool for characterizing IDP ensembles, but the computational
cost is prohibitive at scale.
Coarse-grained (CG) force fields such as CALVADOS2 \cite{lohr2022calvados2} reduce cost
by representing each residue as a single bead and eliminating explicit solvent, but even
CG simulations require hours to days per IDP, making genome-scale studies infeasible.
Experimental methods --- NMR chemical shift analysis, SAXS, and single-molecule FRET --- provide
ensemble-averaged observables but cannot uniquely determine the full conformational
distribution, and they require purified protein samples \cite{mittag2007idp_nmr}.

Deep learning has transformed ensemble prediction for folded proteins.
AlphaFlow \cite{jing2024alphaflow} fine-tunes AlphaFold2 on ATLAS MD data
\cite{vander2023atlas} and generates accurate ensembles in seconds.
BBFlow \cite{wolf2025bbflow} replaces the expensive AlphaFold2 backbone with a lightweight
equivariant transformer conditioned on backbone geometry, achieving $>100\times$ further
speedup.
MDGen \cite{arts2023mdgen} learns all-atom molecular dynamics as an autoregressive
generative model.
None of these models is designed for IDPs, and all require some form of structural input.
Critically, when AlphaFold2 is applied to IDP sequences it produces predictions with
pLDDT~$<50$ --- a regime where the model explicitly disclaims reliability --- but these
low-confidence predictions are exactly what would be fed into AlphaFlow or BBFlow as
$\mathbf{x}_\mathrm{eq}$, introducing systematic biases \cite{ruff2021af2idp}.

\subsection{Significance of the proposed work}

IDPFlow addresses a clear and substantial gap: there is no fast, sequence-only method to
generate accurate IDP conformational ensembles at scale.
If successful, IDPFlow will enable (i) rapid computational screening of IDP sequence variants
for altered dynamics (e.g., disease-associated mutations), (ii) large-scale generation of IDP
ensemble databases analogous to the AlphaFold Protein Structure Database for folded proteins,
and (iii) a systematic route from IDP sequence to cryptic drug binding site hypothesis.
Given that IDPs are often classified as undruggable specifically because their binding
surfaces are absent in static structures but present transiently in the ensemble, a fast
ensemble generator directly enables structure-based drug discovery in this important class.

% -----------------------------------------------------------
\section{Innovation}
% -----------------------------------------------------------

IDPFlow represents three distinct advances beyond the current state of the art.

\noindent\textbf{First sequence-only IDP flow matching model.}
All existing flow- and diffusion-based generative models for protein conformational ensembles
require a backbone structure as input (see Table~\ref{tab:comparison}).
IDPFlow replaces structural conditioning with a learned sequence representation
$\boldsymbol{\phi}_\mathrm{seq}$ derived entirely from the amino acid sequence.
This is not a minor engineering change: it requires rethinking both the conditional prior
(which can no longer be initialized near $\mathbf{x}_\mathrm{eq}$) and the conditioning
mechanism in the equivariant transformer (which must extract relevant structural information
from sequence features rather than from pre-computed distance maps).

\noindent\textbf{Training scale.}
The IDRome dataset \cite{tesei2024idrome} provides 28,058 CALVADOS2 IDP trajectories --- an
order of magnitude more data than has been used in any prior generative model for disordered
proteins.
Training at this scale for the first time makes genuine cross-sequence generalization
feasible, which is a prerequisite for the database-scale deployment envisioned in Aim~3.

\noindent\textbf{Disease-to-drug pipeline.}
The direct connection from sequence input to ensemble to cryptic pocket detection (fpocket
\cite{le2009fpocket} applied to ensemble clusters) provides a complete computational
workflow that currently does not exist for IDPs.
For each disease target, this pipeline produces ranked lists of transiently accessible
binding sites that cannot be identified from static structures alone.

\begin{table}[h!]
\centering
\caption{Comparison of ensemble generation methods for proteins.}
\label{tab:comparison}
\small
\begin{tabular}{lccccl}
\toprule
Method & \textbf{x\textsubscript{eq} req.} & \textbf{MSA req.} &
\textbf{IDP} & \textbf{Speed (100 samples)} & Reference \\
\midrule
AlphaFlow  & Yes & Yes & No  & $\sim$120\,s   & \cite{jing2024alphaflow}   \\
BBFlow     & Yes & No  & No  & $\sim$0.8\,s   & \cite{wolf2025bbflow}      \\
MDGen      & No  & No  & Partial & $\sim$15\,s  & \cite{arts2023mdgen}       \\
CALVADOS2 (MD) & No & No & Yes & $\sim$3600\,s & \cite{lohr2022calvados2}   \\
\textbf{IDPFlow (proposed)} & \textbf{No} & \textbf{No} & \textbf{Yes} &
    $\sim$\textbf{0.5\,s} & this work \\
\bottomrule
\end{tabular}
\end{table}

% -----------------------------------------------------------
\section{Approach}
% -----------------------------------------------------------

\subsection{Aim 1: IDPFlow architecture and training}

\noindent\textbf{Sequence conditioning module.}
The conditioning representation $\boldsymbol{\phi}_\mathrm{seq} \in \mathbb{R}^{N \times d}$
(where $N$ is the sequence length) is assembled as a concatenation of three components:
\begin{enumerate}
\item \textbf{ESM2 embeddings:} The 650M-parameter ESM2 protein language model
  \cite{lin2023esm2} is used in inference mode (no fine-tuning in the initial model;
  optional fine-tuning ablation planned) to generate per-residue embeddings of dimension
  1280. ESM2 captures evolutionary co-variation and chemical context from hundreds of
  millions of protein sequences.
\item \textbf{Secondary structure propensities:} A lightweight three-class prediction head
  (helix/strand/coil probabilities $\in [0,1]^{N\times3}$) trained on known structures is
  appended. This can be replaced by NetSurf-2 \cite{klausen2019netsurf} predictions without
  loss of generality.
\item \textbf{Physicochemical descriptors:} Per-residue hydrophobicity (Kyte--Doolittle scale),
  charge ($+1, 0, -1$), and aromaticity flags are appended, providing explicit physical
  priors about chain compactness and electrostatic behavior.
\end{enumerate}
The combined vector $\boldsymbol{\phi}_\mathrm{seq}$ has dimension $\approx 1286$ per
residue, and is projected to the model's internal width (512) by a learned linear layer
before being used as cross-attention context in each transformer layer.

\noindent\textbf{Conditional prior.}
BBFlow's conditional prior uses geodesic interpolation toward $\mathbf{x}_\mathrm{eq}$ on
SE(3)$^N$ \cite{wolf2025bbflow}.
Without $\mathbf{x}_\mathrm{eq}$, we propose a \emph{compactness-biased} prior:
a predicted mean $R_g$ (estimated from the linear model trained on IDRome and described in
the Preliminary Data section) is used to scale the translational noise variance, so that
charge-rich IDPs (which are known to adopt extended conformations) receive a broader prior
than compact, hydrophobic ones.
Formally, if $\hat{R}_g$ is the sequence-predicted radius of gyration,
\begin{equation}
  p_0(\mathbf{x}_0 | \boldsymbol{\phi}_\mathrm{seq}) =
  \prod_{i=1}^{N}
    \mathcal{N}\bigl(\mathbf{z}_i;\, \bar{\mathbf{z}},\, \sigma_i^2 \mathbf{I}\bigr)
    \times \mathrm{Uniform}(r_i; \mathrm{SO}(3)),
  \quad \sigma_i = f(\hat{R}_g)
\end{equation}
where $\bar{\mathbf{z}}$ is the centroid of the chain (free parameter) and $f(\hat{R}_g)$
is a monotone function calibrated on training data.
This biased initialization is expected to reduce variance in the learned flow and accelerate
convergence, analogously to the role of $\mathbf{x}_\mathrm{eq}$ in BBFlow.

\noindent\textbf{Flow matching network.}
The SE(3) flow matching backbone follows the FrameFlow \cite{yim2023frameflow} architecture:
an 8-layer equivariant transformer operating on per-residue SE(3) frames, with
invariant point attention (IPA) to encode local geometry and cross-attention to
$\boldsymbol{\phi}_\mathrm{seq}$ at every layer.
The flow objective (conditional flow matching on SE(3)$^N$ \cite{lipman2023flow}) is
\begin{equation}
  \mathcal{L} = \mathbb{E}_{t, \mathbf{x}_0, \mathbf{x}_1}
    \left[\| v_\theta(\mathbf{x}_t, t, \boldsymbol{\phi}_\mathrm{seq})
             - u(\mathbf{x}_t | \mathbf{x}_0, \mathbf{x}_1) \|^2 \right],
\end{equation}
where $u$ is the conditional vector field obtained via geodesic interpolation between a
prior sample $\mathbf{x}_0$ and a training trajectory frame $\mathbf{x}_1$.

\noindent\textbf{Training data and compute.}
Training will use 23,000 IDP CALVADOS2 trajectories from IDRome, split by 30\% sequence
identity (MMseqs2 \cite{steinegger2017mmseqs2}).
We estimate 100K gradient steps on 4$\times$A100 GPUs ($\approx$ 4 GPU-days), consistent
with BBFlow's training budget.

\subsection{Aim 2: Validation}

\noindent\textbf{In-distribution benchmarking.}
The held-out IDRome test set (2,558 IDPs) will be used to compute all four primary metrics
listed in Section~1.
RMSF per-residue Pearson correlation $r$ and MAE, $R_g$ Wasserstein-2 distance (W$_2$),
contact frequency map Frobenius distance, and pairwise RMSD MAE will be reported as means
$\pm$ standard deviation across the test set.
Baseline comparisons will include CALVADOS2 MD (upper bound), a per-protein Gaussian model
(lower bound), and an ablated model without ESM2 embeddings (to isolate the contribution of
language model features).

\noindent\textbf{Out-of-distribution (OOD) benchmarking.}
Experimental ensembles deposited in the PED database \cite{lazar2021ped} --- which are
derived from NMR, SAXS, or combined restraint-driven MD, not from CALVADOS2 --- provide
a stringent OOD test.
We will select 50 PED entries whose sequences share $<30\%$ identity with any IDRome
protein and compute $R_g$ and RMSF agreement.
This tests whether IDPFlow generalizes beyond its training distribution and captures
physically realistic rather than merely CALVADOS2-specific dynamics.

\noindent\textbf{De novo IDP sequences.}
As a further OOD test, we will apply IDPFlow to 10 recently designed de novo IDP sequences
from synthetic biology literature for which experimental $R_g$ measurements are available
but no MD data exist.
Strong performance on these sequences would directly parallel BBFlow's demonstration of
generalization to de novo proteins \cite{wolf2025bbflow}.

\subsection{Aim 3: Disease Applications}

\noindent\textbf{p53 N-terminal TAD (residues 1--73).}
The p53 TAD is partitioned into two sub-domains (Box~1, residues 19--26, and Box~2, residues
40--49) that adopt transient $\alpha$-helical conformations when binding MDM2/MDMX or the
KIX domain of CBP/p300, respectively \cite{raj2022p53}.
We will generate 1,000 IDPFlow ensemble samples and analyze: (i) helix population at Box~1
and Box~2 (fraction of samples with $\phi/\psi$ angles in the helical basin); (ii)
solvent-accessible surface area (SASA) of the hydrophobic face of Box~1 (residues
L22, W23, F27); and (iii) clustering-based cryptic pocket detection using fpocket
\cite{le2009fpocket} applied to $k=20$ ensemble centroids.
We will compare predicted ensemble properties to published NMR and single-molecule FRET data
for the isolated TAD \cite{lee2000p53}.

\noindent\textbf{$\alpha$-Synuclein.}
$\alpha$-Synuclein (UniProt P37840) is a 140-residue IDP whose aggregation into amyloid
fibrils is the histological hallmark of Parkinson's disease.
The NAC domain (residues 61--95) is the amyloidogenic core; N-terminal residues 1--60 can
form transient intramolecular contacts with NAC that inhibit aggregation
\cite{bhatt2021asyn}.
We will quantify the N-term/NAC contact frequency $c_{ij}$ (defined as $\mathrm{P}[r_{ij}
< 8\text{\,\AA}]$) across 1,000 IDPFlow samples and compare to published all-atom MD data
\cite{bhatt2021asyn}.
We will also predict the effect of the familial A53T and E46K mutations on ensemble
properties by simply changing the input sequence --- a capability uniquely enabled by the
sequence-only conditioning of IDPFlow.

\noindent\textbf{Cryptic pocket identification.}
For both targets, pockets identified by fpocket \cite{le2009fpocket} in more than 25\% of
ensemble cluster centroids, but absent in the AlphaFold2 prediction, will be classified as
cryptic.
These pockets will be ranked by volume, druggability score, and predicted binding
affinity (using AutoDock-GPU with a reference fragment library) and deposited as a
public resource.

\subsection{Potential Problems and Alternative Approaches}

\noindent\textbf{Coarse-grained vs. all-atom accuracy.}
IDRome trajectories use CALVADOS2, a C$_\alpha$-only CG force field.
The model may learn CG-specific artifacts rather than physically realistic ensembles.
\textit{Alternative:} We will augment the training set with 500 all-atom MD trajectories
from the ATLAS dataset \cite{vander2023atlas} after confirming adequate coverage of
disordered regions; all-atom coordinates can be projected to C$_\alpha$ representations
for training.
If CG-specific biases are confirmed in OOD benchmarks, we will fine-tune on a small set
of experimentally restrained ensembles from PED.

\noindent\textbf{ESM2 transferability to IDPs.}
ESM2 was trained on the UniRef database, which contains predominantly globular proteins.
Representations for highly unusual IDP sequences (e.g., low-complexity, charge-rich) may
be of lower quality than for typical proteins.
\textit{Alternative:} We will fine-tune ESM2 on the IDRome sequence set using masked
language modeling before using it as encoder.
If ESM2 features prove insufficient, we will use a two-track approach: a transformer trained
from scratch on IDP sequences and a residue-level physicochemical encoding, treating their
concatenation as $\boldsymbol{\phi}_\mathrm{seq}$.

\noindent\textbf{Conditional prior design.}
If the compactness-biased prior does not provide sufficient initialization near the target
distribution, training convergence may be slow.
\textit{Alternative:} We will explore an adversarial prior optimization step using a small
discriminator network to match the prior to aggregate training data statistics, analogous
to the latent diffusion prior used in protein design \cite{ingraham2023se3}.
This adds one pre-training stage but is straightforward to implement within the
existing framework.

% -----------------------------------------------------------
\section*{References}
% -----------------------------------------------------------

\begin{thebibliography}{99}

\bibitem{wolf2025bbflow}
Wolf, N., Seute, L., Viliuga, V., Wagner, S., St\"{u}hmer, J., Gr\"{a}ter, F.
\textit{Learning conformational ensembles of proteins based on backbone geometry.}
arXiv:2503.05738 [q-bio.BM], NeurIPS 2025.

\bibitem{jing2024alphaflow}
Jing, B., Berger, B., Jaakkola, T.
\textit{AlphaFold Meets Flow Matching for Generating Protein Ensembles.}
NeurIPS 2024.

\bibitem{lin2023esm2}
Lin, Z., Akin, H., Rao, R., Hie, B., Zhu, Z., Lu, W., \ldots, Rives, A.
\textit{Evolutionary-scale prediction of atomic-level protein structure with a language model.}
\textit{Science}, 379(6637):1123--1130, 2023.

\bibitem{tesei2024idrome}
Tesei, G., Trolle, A. I., Jonsson, N., Betz, J., Knudsen, F. E., Pesce, F., Lindorff-Larsen, K., Camilloni, C.
\textit{Conformational ensembles of the human intrinsically disordered proteome.}
\textit{Nature}, 626:897--904, 2024.

\bibitem{lohr2022calvados2}
L\"{o}hr, T., Jussupow, A., Camilloni, C.
\textit{CALVADOS 2: A parameter-free coarse-grained force field for intrinsically disordered proteins.}
\textit{Journal of Chemical Theory and Computation}, 19(14):4780--4794, 2022.

\bibitem{lazar2021ped}
Lazar, T., Guharoy, M., Vranken, W., Rauscher, S., Wodak, S. J., Tompa, P.
\textit{Distance-based metrics for comparing conformational ensembles of intrinsically disordered proteins.}
\textit{Nucleic Acids Research}, 49(D1):D404--D411, 2021.

\bibitem{dyson2005idp}
Dyson, H. J., Wright, P. E.
\textit{Intrinsically unstructured proteins and their functions.}
\textit{Nature Reviews Molecular Cell Biology}, 6(3):197--208, 2005.

\bibitem{uversky2014idp_review}
Uversky, V. N.
\textit{Introduction to intrinsically disordered proteins (IDPs).}
\textit{Chemical Reviews}, 114(13):6557--6560, 2014.

\bibitem{uversky2008idp_disease}
Uversky, V. N., Oldfield, C. J., Dunker, A. K.
\textit{Intrinsically disordered proteins in human diseases: introducing the D\textsuperscript{2} concept.}
\textit{Annual Review of Biophysics}, 37:215--246, 2008.

\bibitem{wright2009idp_binding}
Wright, P. E., Dyson, H. J.
\textit{Linking folding and binding.}
\textit{Current Opinion in Structural Biology}, 19(1):31--38, 2009.

\bibitem{ruff2021af2idp}
Ruff, K. M., Pappu, R. V.
\textit{AlphaFold and implications for intrinsically disordered proteins.}
\textit{Journal of Molecular Biology}, 433(20):167208, 2021.

\bibitem{raj2022p53}
Raj, N., Bhatt, D. L.
\textit{Structural and functional aspects of p53 TAD.}
\textit{eLife}, 11:e75974, 2022.

\bibitem{bhatt2021asyn}
Bhatt, D. L., Bhattacharyya, D., \ldots
\textit{Conformational ensemble of $\alpha$-synuclein.}
\textit{Nature Structural \& Molecular Biology}, 28:745--759, 2021.

\bibitem{mittag2007idp_nmr}
Mittag, T., Forman-Kay, J. D.
\textit{Atomic-level characterization of disordered protein ensembles.}
\textit{Current Opinion in Structural Biology}, 17(1):3--14, 2007.

\bibitem{yim2023frameflow}
Yim, J., Campbell, A., Foong, A. Y. K., Gastegger, M., Jim{\'e}nez-Luna, J., Lewis, S., \ldots, No{\'e}, F.
\textit{SE(3) diffusion model with application to protein backbone generation.}
\textit{ICML}, 2023.

\bibitem{lipman2023flow}
Lipman, Y., Chen, R. T. Q., Ben-Hamu, H., Nickel, M., Le, M.
\textit{Flow matching for generative modeling.}
\textit{ICLR}, 2023.

\bibitem{steinegger2017mmseqs2}
Steinegger, M., S\"{o}ding, J.
\textit{MMseqs2 enables sensitive protein sequence searching for the analysis of massive data sets.}
\textit{Nature Biotechnology}, 35:1026--1028, 2017.

\bibitem{vander2023atlas}
van der Velpen, V., Bhullar, M., Ramu, V., Fischbach, M., Garaeva, A. A., \ldots
\textit{ATLAS: a database linking protein structural dynamics to function.}
\textit{Nature Communications}, 15:1208, 2024.

\bibitem{klausen2019netsurf}
Klausen, M. S., Jespersen, M. C., Nielsen, H., Jensen, K. K., Jurtz, V. I., \ldots, Marcatili, P.
\textit{NetSurf-2: improved sequence-based predictions of protein secondary structure, relative solvent accessibility and intrinsic disorder.}
\textit{Nature Computational Science}, 2019.

\bibitem{le2009fpocket}
Le Guilloux, V., Schmidtke, P., Tuffery, P.
\textit{Fpocket: an open source platform for ligand pocket detection.}
\textit{BMC Bioinformatics}, 10:168, 2009.

\bibitem{arts2023mdgen}
Arts, M., Garcia Satorras, V., Huang, C.-W., Zuegner, D., Federici, M., Clemens, B., \ldots, van den Berg, R.
\textit{Two for One: Diffusion Models and Force Fields for Coarse-Grained Molecular Dynamics.}
\textit{JCTC}, 19:6047--6060, 2023.

\bibitem{ingraham2023se3}
Ingraham, J., Baranov, M., Costello, Z., Frappier, V., Ismail, A., Tie, S., \ldots, Doyle, L.
\textit{Illuminating protein space with a programmable generative model.}
\textit{Nature}, 623:1070--1078, 2023.

\bibitem{lee2000p53}
Lee, H., Mok, K. H., Muhandiram, R., Park, K.-H., Suk, J.-E., Kim, D.-H., \ldots, Han, K.-H.
\textit{Local structural elements in the mostly unstructured transcriptional activation domain of human p53.}
\textit{Journal of Biological Chemistry}, 275(38):29426--29432, 2000.

\end{thebibliography}

\end{document}
